\begin{enumerate}
\section*{ACADEMIA NOSTRADAMUS SEMESTRAL 2022 - I - semana\_n\_12\_-\_semejanza\_de\_triangulos}
\item[\textbf{1.}] [[curso=Geometria]] [[tema=semejanza de triangulos]] En un triángulo acutángulo $ABC$, se trazan las alturas $\overline{AM}$ y $\overline{BN}$, las cuales se intersectan en $O$. Si $AN=5$, $ON=4$ y $OB=6$, calcule $NC$.£A)$7$æB)$7,5$æC)$8$£D)$8,5$ææE)$9$£ [[clave=-]] [[Estado=sin_revisar]]
\item[\textbf{2.}] [[curso=Geometria]] [[tema=semejanza de triangulos]] Se tiene un triángulo rectángulo $ABC$, recto en $B$. Se ubican $M$, $N$ y $P$ sobre $\overline{AB}$, $\overline{BC}$ y $\overline{AC}$, respectivamente, tal que $MBNP$ es un cuadrado. Si $AM=1$ y $NC=4$, calcule $BP$.£A)$2$æB)$2\sqrt{2}$æC)$2\sqrt{3}$£D)$3$ææE)$4$£ [[clave=-]] [[Estado=sin_revisar]]
\item[\textbf{3.}] [[curso=Geometria]] [[tema=semejanza de triangulos]] En un triángulo rectángulo $ABC$, recto en $B$, se traza una recta $\mathcal{L}$ no secante al triángulo y se traza la bisectriz interior $\overline{BD}$, $m\sphericalangle BCA=37^\circ$. Si $C$ dista de $\mathcal{L}$ $42$, ¿cuánto dista $D$ de $\mathcal{L}$?£A)$28$æB)$27$æC)$18$£D)$12$ææE)$9$£ [[clave=-]] [[Estado=sin_revisar]]
\item[\textbf{4.}] [[curso=Geometria]] [[tema=semejanza de triangulos]] Se tiene un triángulo rectángulo $ABC$ recto en $B$, en el cual se traza la ceviana interior $\overline{AD}$, de modo que $CD=2(BD)$ y $m\sphericalangle DCA=m\sphericalangle DAB$. Calcule $m\sphericalangle BCA$.£A)$37^\circ$æB)$53^\circ$æC)$60^\circ$£D)$45^\circ$ææE)$30^\circ$£ [[clave=-]] [[Estado=sin_revisar]]
\item[\textbf{5.}] [[curso=Geometria]] [[tema=semejanza de triangulos]] Se da un triángulo $ABC$ cuyos lados $\overline{AB}$ y $\overline{BC}$ miden $8$ y $6$, respectivamente. Sobre $\overline{AB}$ se toma el punto $D$, tal que $m\sphericalangle BAC=m\sphericalangle BCD$, entonces $AD$ es:£A)$3,5$æB)$4$æC)$4,5$£D)$5$ææE)$5,5$£ [[clave=-]] [[Estado=sin_revisar]]
\item[\textbf{6.}] [[curso=Geometria]] [[tema=semejanza de triangulos]] En un triángulo $ABC$, sobre la prolongación de $\overline{AC}$ se toma el punto $D$ de tal forma que $4(m\sphericalangle BAC)=m\sphericalangle CDB$. Si $5(m\sphericalangle BAC)=m\sphericalangle ACB$, $BD=\dfrac{10}{\sqrt{3}}$ y $CD=\left(\dfrac{20}{\sqrt{3}}-10\right)$, halle $AC$.£A)$10\sqrt{3}$æB)$20$æC)$12\sqrt{3}$£D)$22$ææE)$13\sqrt{3}$£ [[clave=-]] [[Estado=sin_revisar]]
\item[\textbf{7.}] [[curso=Geometria]] [[tema=semejanza de triangulos]] Sea $ABCD$ un cuadrilátero inscrito en una circunferencia de radio $r$ y circunscrito a una circunferencia de radio $R$. Si $\overline{BD}$ interseca a $\overline{AC}$ en $I$, $3(BI)=AI$ y $AB+CD=a$, calcule $BC$.£A)$\dfrac{a}{2}$æB)$\dfrac{a}{4}$æC)$\dfrac{a}{3}$£D)$\dfrac{a}{5}$ææE)$\dfrac{a}{6}$£ [[clave=-]] [[Estado=sin_revisar]]
\item[\textbf{8.}] [[curso=Geometria]] [[tema=semejanza de triangulos]] En un triángulo $ABC$, se traza la ceviana interior $\overline{BD}$, tal que $m\sphericalangle ABD=m\sphericalangle ACB$, $AB=6$ y $CD=5$. Calcule $\dfrac{BC}{BD}$.£A)$2$æB)$\dfrac{3}{2}$æC)$\dfrac{4}{3}$£D)$\dfrac{6}{5}$ææE)$\dfrac{5}{3}$£ [[clave=-]] [[Estado=sin_revisar]]
\item[\textbf{9.}] [[curso=Geometria]] [[tema=semejanza de triangulos]] Se tienen $2$ circunferencias tangentes exteriores, tangentes en $M$, cuyos radios miden $a$ y $b$. Halle la distancia de $M$ hacia una de las rectas tangentes comunes exteriores.£A)$\sqrt{ab}$æB)$\sqrt{2ab}$æC)$2\sqrt{ab}$£D)$\dfrac{ab}{a+b}$ææE)$\dfrac{2ab}{a+b}$£ [[clave=-]] [[Estado=sin_revisar]]
\item[\textbf{10.}] [[curso=Geometria]] [[tema=semejanza de triangulos]] Por el vértice $A$ de un triángulo $ABC$ inscrito en una circunferencia, se traza una paralela a la tangente de la circunferencia en el punto $B$; dicha paralela interseca a la prolongación del lado $\overline{BC}$ en el punto $E$. Si $AB=6$ y $BC=2$, calcule $CE$.£A)$6$æB)$8$æC)$2\sqrt{3}$£D)$12$ææE)$16$£ [[clave=-]] [[Estado=sin_revisar]]
\item[\textbf{11.}] [[curso=Geometria]] [[tema=semejanza de triangulos]] En una circunferencia de centro $O$ y radio igual a $2,m$, se ubican los puntos consecutivos $A$, $C$, $B$ y $D$, tal que $\overline{AD}$ es diámetro. Desde $B$ se traza $\overline{BH}$ perpendicular a la prolongación de $\overline{AC}$, $H$ en dicha prolongación. Si $(BH)(AC)=3$, calcule la distancia de $C$ a $\overline{AB}$, $B$ pertenece a la mediatriz de $\overline{OD}$.£A)$\sqrt{3}$æB)$\dfrac{\sqrt{3}}{2}$æC)$\dfrac{\sqrt{3}}{3}$£D)$\dfrac{2\sqrt{3}}{3}$ææE)$2\sqrt{3}$£ [[clave=-]] [[Estado=sin_revisar]]
\item[\textbf{12.}] [[curso=Geometria]] [[tema=semejanza de triangulos]] En un triángulo $ABC$, $AB=6$, $BC=7$ y $AC=8$. Calcule la distancia entre el incentro y el baricentro de $ABC$.£A)$1$æB)$2$æC)$\dfrac{1}{2}$£D)$\dfrac{1}{3}$ææE)$\dfrac{2}{3}$£ [[clave=-]] [[Estado=sin_revisar]]
\item[\textbf{13.}] [[curso=Geometria]] [[tema=semejanza de triangulos]] En un triángulo, las longitudes de sus lados son números enteros consecutivos; además, la medida del mayor ángulo interior es el doble del menor ángulo interior. Halle el perímetro de la región triangular inicial.£A)$17$æB)$18$æC)$19$£D)$15$ææE)$12$£ [[clave=-]] [[Estado=sin_revisar]]
\item[\textbf{14.}] [[curso=Geometria]] [[tema=semejanza de triangulos]] En un triángulo $ABC$, se traza la ceviana interior $\overline{BD}$, tal que $m\sphericalangle ACB=2(m\sphericalangle ABD)$, $AB=6$ y $AD=4$. Halle el semiperímetro de la región $ABC$.£A)$6,5$æB)$7$æC)$7,5$£D)$8$ææE)$8,5$£ [[clave=-]] [[Estado=sin_revisar]]
\item[\textbf{15.}] [[curso=Geometria]] [[tema=semejanza de triangulos]] En un triángulo $ABC$, recto en $B$, se traza la altura $\overline{BH}$; en los triángulos $ABH$ y $BHC$ se trazan las alturas $\overline{HM}$ y $\overline{HN}$, respectivamente; en los triángulos $AMH$ y $HNC$ se trazan las alturas $\overline{MF}$ y $\overline{NG}$, respectivamente; luego, en el triángulo $HNG$ se traza la altura $\overline{GI}$ y en el triángulo $HIG$ se traza la altura $\overline{IE}$. Si $MF=a$ y $GN=b$, calcule $IE$.£A)$\sqrt{ab}$æB)$2\sqrt{ab}$æC)$\dfrac{ab}{a+b}$£D)$\dfrac{2ab}{a+b}$ææE)$\dfrac{a+b}{3}$£ [[clave=-]] [[Estado=sin_revisar]]
\item[\textbf{16.}] [[curso=Geometria]] [[tema=semejanza de triangulos]] En un triángulo $ABC$, la mediatriz de $\overline{AC}$ interseca a la circunferencia circunscrita en $P$, y $\overline{AP}$ interseca a $\overline{BC}$ en $D$, tal que $AD=6$ y $AB=4(BD)$. Halle $PD$.£A)$\dfrac{1}{2}$æB)$1$æC)$2$£D)$\dfrac{3}{2}$ææE)$3$£ [[clave=-]] [[Estado=sin_revisar]]
\item[\textbf{17.}] [[curso=Geometria]] [[tema=semejanza de triangulos]] En un triángulo, las longitudes de los lados están en progresión aritmética. Calcule la razón entre el inradio y la longitud de la altura relativa al lado intermedio.£A)$\dfrac{1}{2}$æB)$\dfrac{1}{3}$æC)$\dfrac{1}{4}$£D)$\dfrac{2}{3}$ææE)$\dfrac{1}{5}$£ [[clave=-]] [[Estado=sin_revisar]]
\item[\textbf{18.}] [[curso=Geometria]] [[tema=semejanza de triangulos]] En un cuadrado $ABCD$, en $\overline{BD}$ y $\overline{CD}$ se ubican los puntos $P$ y $Q$, de manera que $m\sphericalangle APQ=90^\circ$; además, las prolongaciones de $\overline{AP}$ y $\overline{DC}$ se intersecan en $S$ y $\overline{PS}\cap\overline{BC}={M}$. Si $(PS)(PM)=16$, calcule $AQ$.£A)$16$æB)$8$æC)$4$£D)$6$ææE)$4\sqrt{2}$£ [[clave=-]] [[Estado=sin_revisar]]
\item[\textbf{19.}] [[curso=Geometria]] [[tema=semejanza de triangulos]] Dado un triángulo $ABC$, se ubican los puntos $M$ y $N$ en $\overline{AC}$ y $\overline{BC}$, respectivamente, tal que los ángulos $BAC$ y $MNC$ son suplementarios. Si $3(AB)=5(MN)$, $NC=6$ y $AC=15$, calcule la longitud de la proyección de $\overline{AB}$ sobre $\overline{AC}$.£A)$2$æB)$\dfrac{5}{2}$æC)$4$£D)$1$ææE)$\sqrt{3}$£ [[clave=-]] [[Estado=sin_revisar]]
\item[\textbf{20.}] [[curso=Geometria]] [[tema=semejanza de triangulos]] En una circunferencia, por los extremos de una cuerda $\overline{AB}$, se trazan rectas tangentes y en el arco menor $\overparen{AB}$ se ubica el punto $P$; además, por dicho punto se traza la recta tangente que interseca a las otras dos en $R$ y $M$. Si las distancias de los puntos $R$ y $M$ a la cuerda $\overline{AB}$ miden $4$ y $9$, respectivamente, calcule la distancia de $P$ a $\overline{AB}$.£A)$\dfrac{36}{13}$æB)$\dfrac{18}{13}$æC)$\dfrac{72}{13}$£D)$\dfrac{26}{9}$ææE)$\dfrac{13}{2}$£ [[clave=-]] [[Estado=sin_revisar]]
\item[\textbf{21.}] [[curso=Geometria]] [[tema=semejanza de triangulos]] En el lado $\overline{AC}$ de un triángulo $ABC$, se ubica el punto $M$, $AM\neq MC$; la circunferencia circunscrita a $BMC$ interseca a $\overline{AB}$ en $P$ y la circunferencia circunscrita a $BMA$ interseca a $\overline{BC}$ en $Q$. Si $(AM)\cdot(MC)=K$, calcule $\sqrt{(PM)(MQ)}$.£A)$K$æB)$\sqrt{K}$æC)$2K$£D)$2\sqrt{K}$ææE)$\sqrt{2K}$£ [[clave=-]] [[Estado=sin_revisar]]
\item[\textbf{22.}] [[curso=Geometria]] [[tema=semejanza de triangulos]] $ABCD$ es un cuadrilátero inscrito en una circunferencia de radio $r$ y circunscrito a una circunferencia de radio $R$. Si $\overline{BD}$ interseca $\overline{AC}$ en $I$, $3(BI)=AI$ y $AB+CD=a,cm$, $a>0$. Calcule la longitud, en $cm$, de $\overline{BC}$.£A)$\dfrac{a}{2}$æB)$\dfrac{a}{3}$æC)$\dfrac{a}{4}$£D)$\dfrac{a}{5}$ææE)$\dfrac{a}{6}$£ [[clave=-]] [[Estado=sin_revisar]]
\item[\textbf{23.}] [[curso=Geometria]] [[tema=semejanza de triangulos]] En un triángulo isósceles $ABC$ de base $\overline{AC}$, calcule la distancia del ortocentro al incentro si $AC=10$ y $AB=13$.£A)$2$æB)$1$æC)$\dfrac{3}{2}$£D)$\dfrac{4}{3}$ææE)$\dfrac{5}{4}$£ [[clave=-]] [[Estado=sin_revisar]]
\item[\textbf{24.}] [[curso=Geometria]] [[tema=semejanza de triangulos]] En un triángulo $ABC$, recto en $B$, se trazan las bisectrices interiores $\overline{AM}$ y $\overline{CN}$, tal que se intersecan en $I$; además, $\dfrac{AN}{CM}=k$. Halle $\dfrac{AH}{HC}$, si $\overline{IH}\perp\overline{AC}$, $H\in\overline{AC}$.£A)$k$æB)$\dfrac{1}{k}$æC)$\dfrac{1}{k+1}$£D)$\dfrac{k}{k+1}$ææE)$\dfrac{k+1}{k}$£ [[clave=-]] [[Estado=sin_revisar]]
\item[\textbf{25.}] [[curso=Geometria]] [[tema=semejanza de triangulos]] Desde un punto exterior a una circunferencia se trazan los segmentos tangentes $\overline{PA}$ y $\overline{PB}$, $A$ y $B$ son puntos de tangencia, y en la prolongación de $\overline{BP}$ se ubica $C$ y se traza $\overline{CH}\perp\overline{AB}$, $H\in\overline{AB}$, y $\overline{BM}$, tal que $\overline{AP}$, $\overline{BM}$ y $\overline{CH}$ son concurrentes en $N$. Si $CP=m$ y $PB=n$, halle $\overline{AE}$, ${E}=\overline{AN}\cap\overline{MH}$.£A)$\dfrac{m}{n}$æB)$\dfrac{n}{m}$æC)$\left(\dfrac{m}{m+n}\right)n$£D)$\left(\dfrac{m+n}{n-m}\right)m$ææE)$\left(\dfrac{m+n}{m}\right)n$£ [[clave=-]] [[Estado=sin_revisar]]
\item[\textbf{26.}] [[curso=Geometria]] [[tema=semejanza de triangulos]] En un triángulo $ABC$, se trazan las cevianas $\overline{BD}$ y $\overline{AE}$, tal que $I$ y $G$ son incentro y baricentro de $ABD$ y $BDC$, respectivamente. Si $m\sphericalangle ABD=2(m\sphericalangle ACB)$, $\overline{IG}\parallel\overline{AC}$, $AB=6$ y $BD=8$, calcule $\dfrac{EC}{EB}$.£A)$1$æB)$2$æC)$3$£D)$4$ææE)$5$£ [[clave=-]] [[Estado=sin_revisar]]
\item[\textbf{27.}] [[curso=Geometria]] [[tema=semejanza de triangulos]] En un trapecio rectángulo $ABCD$, recto en $A$ y $B$, se traza una semicircunferencia de diámetro $\overline{AB}$, la cual es tangente a $\overline{CD}$ en $T$. Si $AD=3$ y $BC=2$, calcule $ET$, $E$: punto de intersección de $\overline{AC}$ y $\overline{BD}$.£A)$\dfrac{2}{3}$æB)$\dfrac{3}{2}$æC)$2$£D)$\dfrac{4}{5}$ææE)$\dfrac{6}{5}$£ [[clave=-]] [[Estado=sin_revisar]]
\item[\textbf{28.}] [[curso=Geometria]] [[tema=semejanza de triangulos]] En un triángulo isósceles, $R$ es radio de la circunferencia inscrita y $r$ el radio de la circunferencia tangente a la primera circunferencia y tangente a los lados laterales. Halle la longitud de la altura relativa a la base.£A)$\dfrac{Rr}{R-r}$æB)$\sqrt{Rr}$æC)$\dfrac{R^2}{R+r}$£D)$\dfrac{2R^2}{R+r}$ææE)$\dfrac{2R^2}{R-r}$£ [[clave=-]] [[Estado=sin_revisar]]
\item[\textbf{29.}] [[curso=Geometria]] [[tema=semejanza de triangulos]] En un cuadrilátero bicéntrico $ABCD$, la circunferencia inscrita es tangente a $\overline{AD}$ en $T$. Si $AB=9$, $CD=10$ y $AT=6$, calcule $BC$.£A)$6,3$æB)$7,3$æC)$6,5$£D)$6,4$ææE)$7,1$£ [[clave=-]] [[Estado=sin_revisar]]
\item[\textbf{30.}] [[curso=Geometria]] [[tema=semejanza de triangulos]] En un triángulo obtusángulo $ABC$ (obtuso en $A$), en la prolongación de $\overline{BA}$ se ubica un punto $T$, luego se traza la ceviana interior $\overline{AM}$, tal que $\overline{AC}$ es bisectriz del ángulo $MAT$ y la mediatriz de $\overline{AC}$ interseca a $\overline{MC}$ en $Q$. Calcule $QC$, si $BM=3$ y $MQ=2$.£A)$5$æB)$6$æC)$5$£D)$6$ææE)$10$£ [[clave=-]] [[Estado=sin_revisar]]
\end{enumerate}
